\clearpage\fancyhead[L]{\small COMET.}\hypertarget{native-387-lower}{}\label{native-387-lower}
\begin{center}\textbf{COMET.}\par\vspace{4pt}\hrule\vspace{4pt}\fontsize{10}{12}\selectfont Göttingische gelehrte Anzeigen. Stück 55. Seite 537\;..\;538. 1815 April 8.\par\vspace{4pt}\hrule\vspace{4pt}\end{center}

\par\indent Vom 13.\ bis 31.\ März waren hier nur vier heitere Abende, wo der neue Comet beobachtet werden konnte. Folgendes ist daher die ganze auf der hiesigen Sternwarte bisher erhaltene Ausbeute von Beobachtungen:

\begin{center}\fontsize{9}{11}\selectfont\setlength{\tabcolsep}{3pt}\renewcommand{\arraystretch}{1.15}\begin{tabular}{l|ccc|ccc|ccc}
1815 & \multicolumn{3}{c|}{Mittlere Zeit} & \multicolumn{3}{c|}{\shortstack{Scheinb. gerade\\Aufsteigung}} & \multicolumn{3}{c}{\shortstack{Scheinbare\\Abweichung}} \\

\hline
März 20 & $10^{\mathrm u}$ & $33^{\mathrm m}$ & $6^{\mathrm s}$ & $54^\circ$ & $7'$ & $1''$ & $39^\circ$ & $7'$ & $47''$ N. \\
21 & 10 & 11 & 37 & 54 & 33 & 21 & 39 & 36 & 57 \\
21 & 10 & 57 & 0 & 54 & 34 & 21 & 39 & 38 & 8 \\
25 & 9 & 46 & 38 & 56 & 28 & 50 & 41 & 38 & 5 \\
30 & 9 & 50 & 57 & 59 & 13 & 3 & 44 & 10 & 27 \\
\end{tabular}\end{center}

\par\indent Die zweite Beobachtung vom 21.\ März ist von Hrn.\ \textsc{Enke}, von dessen Geschicklichkeit in practisch-astronomischen Geschäften diese Blätter schon öfters Beweise gegeben haben. Weitere auswärtige Beobachtungen sind noch nicht eingegangen.

\Needspace{19\baselineskip}\par\indent Die vorläufige genäherte Bestimmung der Bahn dieses Cometen, welche Hr.\ Prof.\ \textsc{Gauss} nach diesen freilich nur dürftigen Beobachtungen gemacht hat, ist folgende:

\begin{center}\fontsize{10}{12}\selectfont\setlength{\tabcolsep}{3pt}\renewcommand{\arraystretch}{1.15}\begin{tabular}{l}
Durchgang durch die Sonnennähe\\
\quad Mittl. Zeit in Göttingen 1815 April 24. $16^{\mathrm u}\;37^{\mathrm m}\;34^{\mathrm s}$\\
\end{tabular}\end{center}
\begin{center}\fontsize{10}{12}\selectfont\setlength{\tabcolsep}{3pt}\renewcommand{\arraystretch}{1.15}\begin{tabular}{lcr}
Länge der Sonnennähe & $.\;.\;.\;.\;.$ & $246^\circ\;7'\;2''$\\
Aufsteigender Knoten & $.\;.\;.\;.\;.\;.$ & $82\;43\;4$\\
Neigung der Bahn & $.\;.\;.\;.\;.\;.\;.$ & $45^\circ\;8'\;55''$\\
Kleinster Abstand von der Sonne & $.\;.$ & $1.24738$\\
\end{tabular}\end{center}
\noindent Bewegung rechtläufig.

\SourceLeaf{388}\par\indent Es lässt sich aus dieser Bestimmung bereits voraussehen, dass der Comet bis gegen Ende Aprils noch etwas an Licht zunehmen wird. Diese Zunahme ist aber sehr gering; er wird dem blossen Auge unsichtbar bleiben, und daher \emph{nur für die Astronomen} interessant sein. Er wird durch den Fahrmann und Luchs in den grossen Bär gehen, und dort, wegen seines immer mehr abnehmenden Lichts, im Juni sich dem Auge entziehen.
